\vspace{0.5em}\noindent\textbf{Self-evaluation}\par

During the development of the \textbf{Internship Application Platform},
I had the opportunity to apply knowledge in full-stack development,
Cloud architecture, AI integration, realtime communication,
containerization, Kubernetes, CI/CD, and observability. This
self-evaluation summarizes my completion level, strengths, limitations,
and improvement direction after the internship period.

\vspace{0.5em}\noindent\textbf{Evaluation Criteria}\par

{
\begingroup
\small
\setlength{\tabcolsep}{3pt}
\renewcommand{\arraystretch}{1.15}
\begin{longtable}{@{}
  >{\raggedright\arraybackslash}p{(\linewidth - 10\tabcolsep) * \real{0.1667}}
  >{\raggedright\arraybackslash}p{(\linewidth - 10\tabcolsep) * \real{0.1667}}
  >{\raggedright\arraybackslash}p{(\linewidth - 10\tabcolsep) * \real{0.1667}}
  >{\raggedright\arraybackslash}p{(\linewidth - 10\tabcolsep) * \real{0.1667}}
  >{\raggedright\arraybackslash}p{(\linewidth - 10\tabcolsep) * \real{0.1667}}
  >{\raggedright\arraybackslash}p{(\linewidth - 10\tabcolsep) * \real{0.1667}}@{}}
\toprule\noalign{}
\begin{minipage}[b]{\linewidth}\raggedright
No.
\end{minipage} & \begin{minipage}[b]{\linewidth}\raggedright
Criteria
\end{minipage} & \begin{minipage}[b]{\linewidth}\raggedright
Description
\end{minipage} & \begin{minipage}[b]{\linewidth}\raggedright
Good
\end{minipage} & \begin{minipage}[b]{\linewidth}\raggedright
Fair
\end{minipage} & \begin{minipage}[b]{\linewidth}\raggedright
Average
\end{minipage} \\
\midrule\noalign{}
\endhead
\bottomrule\noalign{}
\endlastfoot
1 & \textbf{Professional knowledge \& skills} & Understanding of the
field, applying knowledge in practice, proficiency with tools, and work
quality & \checkmark & \(\square\) & \(\square\) \\
2 & \textbf{Ability to learn} & Ability to absorb new knowledge and
learn quickly & \checkmark & \(\square\) & \(\square\) \\
3 & \textbf{Proactiveness} & Taking initiative and seeking out tasks
without waiting for instructions & \checkmark & \(\square\) &
\(\square\) \\
4 & \textbf{Sense of responsibility} & Completing tasks on time and
ensuring quality & \checkmark & \(\square\) & \(\square\) \\
5 & \textbf{Discipline} & Adhering to schedules, rules, and work
processes & \(\square\) & \checkmark & \(\square\) \\
6 & \textbf{Progressive mindset} & Willingness to receive feedback and
improve oneself & \checkmark & \(\square\) & \(\square\) \\
7 & \textbf{Communication} & Presenting ideas and reporting work clearly
& \(\square\) & \checkmark & \(\square\) \\
8 & \textbf{Teamwork} & Working effectively with colleagues and
participating in teams & \checkmark & \(\square\) & \(\square\) \\
9 & \textbf{Professional conduct} & Respecting colleagues, partners, and
the work environment & \checkmark & \(\square\) & \(\square\) \\
10 & \textbf{Problem-solving skills} & Identifying problems, proposing
solutions, and showing creativity & \checkmark & \(\square\) &
\(\square\) \\
11 & \textbf{Contribution to project/team} & Work effectiveness,
technical contribution, and recognition from the team & \checkmark &
\(\square\) & \(\square\) \\
12 & \textbf{Overall} & General evaluation of the entire internship
period & \checkmark & \(\square\) & \(\square\) \\
\end{longtable}
\endgroup
}

\vspace{0.5em}\noindent\textbf{Strengths}\par

\begin{itemize}
\tightlist
\item
  I can approach problems from a system perspective instead of focusing
  only on isolated features. I considered authentication, storage,
  realtime chat, AI processing, deployment, and observability together.
\item
  I proactively learned AWS services and Cloud-native tools to
  understand how an application can operate in a production-oriented
  environment.
\item
  I connected academic knowledge with practical project requirements,
  such as using S3 presigned URLs for document security, Redis adapter
  for realtime chat, and workers for long-running tasks.
\item
  I improved my technical writing skills, bilingual report structure,
  and ability to follow the FCAJ report template.
\item
  I was open to feedback and willing to revise the content when
  something was not clear or consistent.
\end{itemize}

\vspace{0.5em}\noindent\textbf{Areas for Improvement}\par

\begin{itemize}
\tightlist
\item
  I need to improve time estimation for each work package, especially
  parts involving many components such as Kubernetes, observability, and
  CI/CD.
\item
  I need more hands-on practice with production deployment on AWS,
  including networking, IAM, cost control, and real monitoring.
\item
  I need to improve automated testing coverage for auth, applications,
  documents, chat, and AI flows.
\item
  I need to improve UI/UX design skills so that the interface is easier
  to use for both Candidates and HR users.
\item
  I need to keep improving how AI results are explained, so that
  matching scores are not treated as the only recruitment decision.
\end{itemize}

\vspace{0.5em}\noindent\textbf{Future Learning Plan}\par

\begin{itemize}
\tightlist
\item
  Review and practice the AWS Well-Architected Framework, especially
  security, reliability, and cost optimization.
\item
  Continue learning Kubernetes, Helm, Ingress, HPA, PDB, probes, and
  deployment strategies.
\item
  Study observability with Prometheus, Grafana, Loki, Tempo, and
  OpenTelemetry.
\item
  Improve backend skills with FastAPI, SQLAlchemy, Alembic, PostgreSQL
  transactions, and security best practices.
\item
  Practice CI/CD with GitHub Actions, OIDC, container scanning, secret
  scanning, and deployment automation.
\item
  Learn more about AI application design, prompt engineering, schema
  validation, fallback strategy, and human-in-the-loop review.
\end{itemize}

\vspace{0.5em}\noindent\textbf{Career Orientation}\par

After the internship, I would like to continue developing toward
\textbf{Cloud/Backend Engineering} combined with \textbf{AI
integration}. This project helped me understand that a real system does
not only need working features. It also needs security, scalability,
observability, deployment automation, and cost control.

In the future, I want to keep building more practical projects, deploy
them on AWS in a controlled way, and strengthen my system design mindset
so that I can contribute better to Cloud-native projects in a
professional environment.
